\answer{OR26}
令
\[
 \tau=\inf\{t\ge0:B_t^{(1)}=-1\}.
\]
由反射原理，$\Prob(\tau\le t)=2\Phi(-1/\sqrt t)$，其中 $\Phi$ 为标准正态分布函数。因此 $\tau<\infty$ 几乎必然成立，其密度为
\[
 f_\tau(t)=\frac1{\sqrt{2\pi t^3}}e^{-1/(2t)},\qquad t>0.
\]
$\tau$ 只由第一坐标决定，与整个第二坐标过程独立，所以给定 $\tau=t$，有 $B_\tau^{(2)}\sim N(0,t)$。以 $Z=B_\tau^{(2)}$ 记第二坐标的增量。混合正态密度得到
\[
 f_Z(z)=\int_0^\infty\frac1{2\pi t^2}
            e^{-(1+z^2)/(2t)}dt
       =\frac1{\pi(1+z^2)}.
\]
这里被积函数非负，混合与积分交换由 Tonelli 定理保证；最后令 $u=1/t$ 即得积分值。于是 $Z$ 服从标准柯西分布，且
\[
 \boxed{\Prob(1+Z>0)
   =\int_{-1}^\infty\frac{dz}{\pi(1+z^2)}
   =\frac34}.
\]
独立的是停止时间与第二坐标的整个过程；随机时刻取值的无条件分布是上述混合分布，而不是固定方差的正态分布。
